\emph{Step 1 (clause (i)).} The rung field is that of \claimref{synthesis.rung-field-on-the-lift-locus}, with the stored slot action. In $\Cl(3,5)$ take $v_{a} = \tfrac12(e_{a} + f_{a})$, $w_{a} = \tfrac12(e_{a} - f_{a})$ for the rung planes $a = 4, 5$, so $v_{a}^{2} = w_{a}^{2} = 0$, $\{v_{a}, w_{b}\} = \delta_{ab}$, and $I_{4}$ commutes with every $v_{a}$, $w_{a}$ (\claimref{synthesis.linear-rung-field-has-a-neutral-normal-pair-in-the-scalon-sum}(1)). For $x, y$ in $\{v_{4}, v_{5}, w_{4}, w_{5}\}$, $[xy, z] = x\{y, z\} - \{x, z\}y$ for a vector $z$, so $[v_{4}w_{5}, v_{5}] = v_{4}$, $[v_{4}w_{5}, w_{4}] = -w_{5}$, $[v_{5}w_{4}, v_{4}] = v_{5}$, $[v_{5}w_{4}, w_{5}] = -w_{4}$, and $[b_{a}, v_{a}] = v_{a}$, $[b_{a}, w_{a}] = -w_{a}$ (loc. cit. Step 2). Reading the coefficients of $c$ gives: $\tfrac12[I_{4}(v_{4}w_{5} + v_{5}w_{4}), \cdot]$ maps $\alpha \mapsto \tfrac{I_{4}}{2}\sigma_{1}\alpha$, $\beta \mapsto -\tfrac{I_{4}}{2}\sigma_{1}\beta$; $\tfrac12[v_{4}w_{5} - v_{5}w_{4}, \cdot]$ maps both $\alpha$ and $\beta$ by the real matrix $\tfrac{I_{4}}{2}\sigma_{2}$; $\tfrac12[I_{4}(b_{4} \mp b_{5}), \cdot]$ maps $\alpha \mapsto \tfrac{I_{4}}{2}\sigma_{3}\alpha$ or $\tfrac{I_{4}}{2}\alpha$ and $\beta$ by the conjugates. These are the stored $T_{1}$, $T_{2}$, $T_{3}$, $Y$ with the unit $I_{4}$. The coefficients are read back by $\alpha_{a} = \{w_{a}, c\}$, $\beta_{a} = \{v_{a}, c\}$, and $Q = c^{2}$ by the Clifford product. The check builds the four generators and the form $Q$ in the GA engine and verifies them against the stored action, and verifies $\nabla V\cdot t\varphi = 0$ for all four. For the coupling: by \claimref{synthesis.rung-field-coupling-to-its-conjugate-is-the-second-family-matrix}, at a point with only rung-$4$ components $M_{d} = \beta Y + \bar{\alpha}\tilde{Y}$ and $M_{u} = \alpha Y + \bar{\beta}\tilde{Y}$; with $Y_{33}$ neglected the third family has $M_{u} = \bar{\beta}_{4}y$, $M_{d} = \bar{\alpha}_{4}y$, whose $SU(2)$-covariant completion in Weyl fields $q$, $u^{c}$, $d^{c}$ is $y\,\epsilon(q, \bar\beta)u^{c} - y\,(q\cdot\bar\alpha)d^{c}$ (under $X \in SU(2)$, $\bar{\beta} \mapsto X\bar{\beta}$, $\bar{\alpha} \mapsto \bar{X}\bar{\alpha}$, and $\epsilon(Xq, X\bar\beta) = \epsilon(q, \bar\beta)$, $(Xq)\cdot(\bar X\bar\alpha) = q\cdot\bar\alpha$); the hypercharges balance, $\tfrac16 + \tfrac12 - \tfrac23 = 0$ and $\tfrac16 - \tfrac12 + \tfrac13 = 0$. The check verifies the invariance $M_{F}(t\varphi) + \tau^{\mathsf T}M_{F} + M_{F}\tau = 0$ for the four generators exactly, and that at $\alpha = 0$, $\beta = (-1, 0)$ only the top pair has mass. Both couplings equal $\tilde{Y}_{33}$, so $y_{1} = y_{2} = m_{t}/v$ at the vacuum.

\emph{Step 2 (clause (ii)).} The gauge loops are taken in Landau gauge with the indefinite form of the Krein premise \claimref{synthesis.krein-form-premise-indefinite-inner-product-for-gauge-fields} on the gauge field, as in the stored one-loop running; the scalar and Yukawa terms do not use it. For real scalars $\varphi$ with a scale-free quartic $V$, vector masses $M_{V}^{2}(\varphi)_{AB} = g_{A}g_{B}(t_{A}\varphi)\cdot(t_{B}\varphi)$ and Weyl mass matrix $M_{F}(\varphi)$, the one-loop effective potential in Landau gauge is $V + \tfrac{1}{64\pi^{2}}\mathrm{STr}\,M^{4}(\ln M^{2}/\mu^{2} - c)$ (Fact \claimref{synthesis.imported-coleman-weinberg-gildener-weinberg-scalon-mass}(i)), with $\mathrm{STr}\,M^{4} = \operatorname{Tr}V''^{2} + 3\operatorname{Tr}M_{V}^{4} - 2\operatorname{Tr}(M_{F}M_{F}^{\dagger})^{2}$. Its independence of $\mu$ at this order, together with the field's anomalous dimension $\gamma = Y_{2}(S) + 3\sum_{A}g_{A}^{2}t_{A}t_{A}$, $Y_{2}(S)_{ab} = \tfrac12\operatorname{Tr}(Y^{a\dagger}Y^{b} + Y^{b\dagger}Y^{a})$, $Y^{a} = \partial M_{F}/\partial\varphi_{a}$, gives
\[
16\pi^{2}\beta_{V}(\varphi) = \tfrac12\bigl(\operatorname{Tr}V''^{2} + 3\operatorname{Tr}M_{V}^{4} - 2\operatorname{Tr}(M_{F}M_{F}^{\dagger})^{2}\bigr) + \varphi\cdot\gamma\nabla V,
\]
the one-loop part of the general formula \cite{MachacekVaughn1985, LuoWangXiao2003} written as a polynomial (cited from memory). The check confirms the normalisation on one doublet ($V = \lambda|H|^{4}$ gives $24\lambda^{2} - 3\lambda(3g^{2} + g'^{2}) + \tfrac38(2g^{4} + (g^{2} + g'^{2})^{2}) + 12\lambda y_{t}^{2} - 6y_{t}^{4}$) and on $O(N)$, $N = 1, 8$ ($(N + 8)/3$ times $\lambda^{2}$ for $V = \tfrac{\lambda}{24}(\varphi^{2})^{2}$ \cite{KleinertSchulteFrohlinde2001}). Since $\beta_{V}$ is $U(2)$-invariant and quartic it is a combination of the ten invariants; the check computes the fourth-derivative tensor of $V$ by polarisation, $\beta_{V}$ at sampled points and couplings, and compares with $\sum_{i}\dot k_{i}P_{i}(\varphi)$ for the closed forms. The general forms are those of clause (ii) plus, with $\kappa = 12(k_{4}^{2} + k_{5}^{2})$ and $\kappa' = 12(k_{6}^{2} + k_{7}^{2})$: $\kappa$ in $16\pi^{2}\dot k_{1}$, $\kappa'$ in $16\pi^{2}\dot k_{2}$, $4(k_{4}^{2} + k_{5}^{2} + k_{6}^{2} + k_{7}^{2}) + 16(k_{4}k_{6} + k_{5}k_{7})$ in $16\pi^{2}\dot k_{3}$, $10(k_{4}^{2} - k_{5}^{2} + k_{6}^{2} - k_{7}^{2}) + 4(k_{4}k_{6} - k_{5}k_{7})$ in $16\pi^{2}\dot k_{8}$, $20(k_{4}k_{5} + k_{6}k_{7}) + 4(k_{4}k_{7} + k_{5}k_{6})$ in $16\pi^{2}\dot k_{9}$, $10(k_{4}^{2} + k_{5}^{2} + k_{6}^{2} + k_{7}^{2}) + 4(k_{4}k_{6} + k_{5}k_{7})$ in $16\pi^{2}\dot k_{10}$, and
\begin{align*}
16\pi^{2}\dot k_{4} &= (96k_{1} + 24k_{3} + 32k_{10} + 40k_{8})k_{4} + 40k_{9}k_{5} + (24k_{3} + 16k_{10} + 8k_{8})k_{6} + 8k_{9}k_{7} - wk_{4} + (18y_{1}^{2} + 6y_{2}^{2})k_{4},\\
16\pi^{2}\dot k_{5} &= (96k_{1} + 24k_{3} + 32k_{10} - 40k_{8})k_{5} + 40k_{9}k_{4} + (24k_{3} + 16k_{10} - 8k_{8})k_{7} + 8k_{9}k_{6} - wk_{5} + (18y_{1}^{2} + 6y_{2}^{2})k_{5},
\end{align*}
and $\dot k_{6}$, $\dot k_{7}$ the same with $(k_{1}, k_{4}, k_{5}, y_{1}, y_{2}) \leftrightarrow (k_{2}, k_{6}, k_{7}, y_{2}, y_{1})$. Every term of $\dot k_{4}, \ldots, \dot k_{7}$ contains one of $k_{4}, \ldots, k_{7}$, so the branch $k_{4} = \cdots = k_{7} = 0$ is preserved; it is the branch of the stored physical ray when also $k_{6} = k_{7} = 0$ there is required (\claimref{synthesis.linear-rung-field-has-a-neutral-normal-pair-in-the-scalon-sum}(3)). The coupling is invariant under $\alpha \to -\alpha$, $d^{c} \to -d^{c}$, which is why no $k_{4}, \ldots, k_{7}$ is generated; the neglected $Y_{33}$ breaks it at order $y_{t}^{3}y_{b}$.

\emph{Step 3 (dictionary).} With canonical real components $|\beta|^{2} = v^{2}$, so the doublets of the standard normalisation are $\Phi = \beta/\sqrt2$, $\alpha/\sqrt2$ (with the conjugations of Step 1), and $|\Phi_{1}^{\dagger}\Phi_{2}| = \tfrac12|Q|$: $\tfrac{\lambda_{1}}{2}|\Phi_{1}|^{4} = k_{1}A^{2}$ gives $\lambda_{1} = 8k_{1}$, and likewise $\lambda_{2} = 8k_{2}$, $\lambda_{3} = 4k_{3}$, $\lambda_{4} = 4k_{10}$, $\operatorname{Re}\lambda_{5} = 4k_{8}$, $|\operatorname{Im}\lambda_{5}| = 4|k_{9}|$; the top coupling of the canonical doublet is $\sqrt2 y_{2} = y_{t}$. The check substitutes these into the standard one-loop equations of the two-doublet model with a discrete symmetry and type-II couplings (Fact \claimref{synthesis.imported-general-quartic-running-and-two-doublet-renormalisation-group-equations}(ii): $16\pi^{2}\dot\lambda_{1} = 12\lambda_{1}^{2} + 4\lambda_{3}^{2} + 4\lambda_{3}\lambda_{4} + 2\lambda_{4}^{2} + 2|\lambda_{5}|^{2} + \tfrac34(3g^{4} + g'^{4} + 2g^{2}g'^{2}) - 3\lambda_{1}(3g^{2} + g'^{2}) + 12\lambda_{1}y_{b}^{2} - 12y_{b}^{4}$, and the forms for $\lambda_{2}$ to $\lambda_{5}$) and finds them equal to the closed forms term by term at sampled couplings. The masses of \claimref{synthesis.linear-rung-field-has-a-neutral-normal-pair-in-the-scalon-sum}(3), $m_{\pm}^{2} = 2k_{3}v^{2}$, $m_{H,A}^{2} = m_{\pm}^{2} + 2v^{2}(k_{10} \pm r)$, are the familiar $\tfrac12\lambda_{3}v^{2}$ and $\tfrac12(\lambda_{3} + \lambda_{4} \pm |\lambda_{5}|)v^{2}$ of a doublet without a vacuum value.

\emph{Step 4 (clause (iii)).} For $V = \tfrac{1}{24}\lambda_{abcd}\varphi_{a}\varphi_{b}\varphi_{c}\varphi_{d}$ the pure-quartic two-loop term of the general formula, $-\tfrac14\sum_{\mathrm{perms}}\lambda_{abef}\lambda_{cegh}\lambda_{dfgh} + \sum_{\mathrm{legs}}\gamma^{(2)}\lambda$ with $\gamma^{(2)}_{ab} = \tfrac{1}{12}\lambda_{acde}\lambda_{bcde}$ (Fact \claimref{synthesis.imported-general-quartic-running-and-two-doublet-renormalisation-group-equations}(i)), contracted with $\varphi^{4}/24$ is $-\tfrac12V''_{ef}V'''_{egh}V'''_{fgh} + \tfrac{1}{12}(\lambda\lambda)_{ab}\varphi_{a}\partial_{b}V$, since $\lambda_{abef}\varphi_{a}\varphi_{b} = 2V''_{ef}$ and $\lambda_{cegh}\varphi_{c} = V'''_{egh}$. One real scalar gives $-\tfrac14\cdot24 + \tfrac13 = -\tfrac{17}{3}$, and $O(N)$ gives $-(3N + 14)/3$ \cite{KleinertSchulteFrohlinde2001}; the check confirms both for $N = 1$ and $N = 8$, then compares the closed forms of clause (iii) with the field computation at three points; $\dot k_{4}^{(2)}, \ldots, \dot k_{7}^{(2)}$ vanish on the branch. For orientation (computed): at the symmetric point $k_{1} = k_{2} = 0$, $k_{3} = \tfrac65$, $k_{10} = s = 0$ the two-loop term of $\dot k_{3}$ is $-192k_{3}^{3}/(16\pi^{2})^{2}$ against $16k_{3}^{2}/(16\pi^{2})$, a ratio of about $-0.09$.

\emph{Step 5 (clause (iv)).} On the ray $\varphi = (0, 0, 0, 0, -v, 0, 0, 0)$ with $k_{2} = k_{6} = k_{7} = 0$, $\nabla V = 0$ (every term is at least quadratic in $\alpha$ or proportional to $k_{2}$), so $16\pi^{2}\beta_{V} = \tfrac12\mathrm{STr}\,M^{4}$ there and $\beta_{V} = \dot k_{2}v^{4}$. Directly from clause (ii) at $k_{2} = 0$: $2v^{4}(8k_{3}^{2} + 8k_{3}k_{10} + 4k_{10}^{2} + 4s) = 8k_{3}^{2}v^{4} + 8v^{4}((k_{3} + k_{10})^{2} + s) = 2m_{\pm}^{4} + m_{H}^{4} + m_{A}^{4}$, $2v^{4}G = \tfrac{1}{16}(6g^{4} + 3(g^{2} + g'^{2})^{2})v^{4} = 6M_{W}^{4} + 3M_{Z}^{4}$, and $2v^{4}(-6y_{2}^{4}) = -12m_{t}^{4}$. Along the ray $V_{1\text{-loop}} = (k_{2}(\mu) + B\ln(\rho^{2}/\mu^{2}))\rho^{4}$ up to a constant times $\rho^{4}$, and $\mu$-independence gives $B = \tfrac12\dot k_{2}$; with Fact \claimref{synthesis.imported-coleman-weinberg-gildener-weinberg-scalon-mass}(ii), $m_{h}^{2} = 8Bv^{2} = 4v^{2}\dot k_{2}$, which is \claimref{synthesis.linear-rung-field-has-a-neutral-normal-pair-in-the-scalon-sum}(4) and, at its standard reading, Fact \claimref{synthesis.imported-two-doublet-scalon-sum-rule-and-lep-neutrino-count}(i). The check verifies the identity at three samples both from the closed form and from the field computation on the ray, and that dropping $m_{H}^{4} + m_{A}^{4}$ breaks it. For the self-coupling, with $L = \ln(\rho^{2}/v^{2})$ and $dL/d\rho = 2/\rho$, $V = B\rho^{4}(L - \tfrac12)$ has $V' = 4B\rho^{3}L$, $V'' = 12B\rho^{2}L + 8B\rho^{2}$, $V''' = 24B\rho L + 40B\rho$, so at $\rho = v$: $V' = 0$, $V'' = 8Bv^{2} = m_{h}^{2}$, $V''' = 40Bv = \tfrac53(3m_{h}^{2}/v)$, the tree value for the same mass being $3m_{h}^{2}/v$ (standard for Gildener-Weinberg potentials). By homogeneity the scalon couples to each heavy pair with $\partial m_{X}^{2}/\partial\rho = 2m_{X}^{2}/v$, the unit weight used in \claimref{synthesis.linear-rung-field-has-a-neutral-normal-pair-in-the-scalon-sum}(5).

\emph{Step 6 (clause (v), computed).} Initial values at $\mu_{0} = m_{t}$ (from memory, $\overline{\mathrm{MS}}$): $g = 0.648$, $g' = 0.358$, $g_{3} = 1.164$, $y_{t} = 0.936$, so $y_{1} = y_{2} = 0.662$; the gauge couplings run at one loop with the stored two-doublet coefficients $(\tfrac{21}{5}, -3, -7)$ of \claimref{synthesis.one-loop-running-coefficients-from-spacetime-algebra-spin-traces} (gauge loops, so under the Krein premise there too), that is $(7, -3, -7)$ for $(g', g, g_{3})$ with $g'$ in the non-GUT normalisation, and $y_{1} = y_{2}$ with $16\pi^{2}\dot y_{2} = y_{2}(y_{1}^{2} + 9y_{2}^{2} - 8g_{3}^{2} - \tfrac94g^{2} - \tfrac{17}{12}g'^{2})$ (the Yukawa part from the general one-loop formula applied to the coupling of Step 1, computed; its $y_{1} = 0$ limit is the standard $\tfrac92y_{t}^{3}$), with the hypercharge term of the $d^{c}$ coupling, $-\tfrac{5}{12}g'^{2}$, averaged in (negligible here). The quartics start at $k_{3} = m_{\pm}^{2}/(2v^{2})$, $k_{10} = (m_{H}^{2} + m_{A}^{2} - 2m_{\pm}^{2})/(4v^{2})$, $|k_{8}| = |m_{H}^{2} - m_{A}^{2}|/(4v^{2})$, $k_{9} = 0$, $k_{1} = k_{2} = 0$. Fourth-order Runge-Kutta in $t$, step $0.002$ to $0.01$, over a grid of points of the surface (both ends of the sum, $m_{\pm}$ from $80$ to $450$ GeV, the share of $m_{A}^{4}$ from $0$ to $\tfrac12$, $m_{A} > M_{Z}/2$): the scale at which $\max_{i}|\lambda_{i}| = 4\pi$ is between $0.80$ TeV ($m_{\pm} = 80$, $m_{H} = 528$, $m_{A} = 305$ GeV) and $4.0$ TeV ($m_{\pm} = 320$, $m_{H} = m_{A} = 418$ GeV) at one loop, between $1.08$ and $15.1$ TeV with the two-loop pure-quartic terms, which at the one-loop crossing are already $-0.67$ (degenerate end) to $-0.31$ (light charged end) times the one-loop term of the largest coupling, so this shift is not a controlled estimate and is not part of the claim; the first coupling to reach $4\pi$ is $\lambda_{3}$ near the degenerate point and $\lambda_{4}$ for a light charged pair. Starting at $\mu_{0} = v$ instead raises the scales by about $\tfrac{v}{m_{t}} \approx 1.4$; $k_{1}(\mu_{0}) = \tfrac12$ lowers the one-loop scale to about $0.6$ to $0.9$ TeV. Since $k_{1}$ and $k_{2}$ start equal and $y_{1} = y_{2}$, they run equally on the branch. The check reproduces the two endpoints in exact arithmetic at one loop: the closed forms of clause (ii) with $y_{1} = y_{2} = y$, the gauge and top flow above, fourth-order Runge-Kutta in $t$ with step $\tfrac{1}{25}$ and states rounded to $10^{-12}$, run with $\pi$ replaced by each of $\tfrac{333}{106} < \pi < \tfrac{355}{113}$; the first step with $\max_{i}|\lambda_{i}| \ge 4\pi$ is step $39$ and step $79$, which by rational bounds on $e^{t}$ puts the crossing between $0.75$ and $0.85$ TeV for $(80, 528, 305)$ GeV and between $3.8$ and $4.2$ TeV for $(320, 418, 418)$ GeV (computed; the rounding and the step size are not error-controlled, so this is a reproducible bracket of the computed flow, not a proof).

\emph{Reading.} \claimref{synthesis.covariant-rung-coupling-gives-the-normal-pair-tree-level-b-meson-mixing-that-excludes-the-transmutation-branch} excludes the same branch under the same coupling reading by tree-level $B$-meson mixing. Standard: the one-loop two-doublet equations, the general two-loop formula, the Gildener-Weinberg relation between the flat-direction running and the scalon mass, the $\tfrac53$ self-coupling. New: the equations in the series' invariants with the series' generators and coupling reading (both doublets at top strength), the two-loop pure-quartic terms in those invariants, the check that the stored sum with the normal pair is that running, and the strong-coupling scale of the transmutation branch, which says the elementary rung field on that branch is a perturbative theory only up to a few TeV. The running is regulator independent at this order (only $\mu$ enters). Transport-convenient.
